\documentclass[00-main.tex]{subfiles}

\begin{document}

\section{Introduction}

\subsection{Perceptual Decision Making in Naturalistic Settings}
Perceptual decision making is the process of making decisions based on sensory information. While often studied in highly controlled experiments, in naturalistic settings these decisions are rarely stable; rather, they emerge from the interplay between multiple interleaved decision states that are richly structured in time \citep{roy_efficient_2018, roy_extracting_2021, ashwood_mice_2022,duffy_disentangling_2025, urai_structure_2025, danskin_exponential_2023, perquin_temporal_2023, cochrane_multiple_2023}. The dominant theoretical models of these two-alternative forced-choice (2AFC) decisions are sequential sampling models, which postulate that decisions are made upon the sequential analysis of evidence \citep{shadlen_speed_2006, asadpour_rodents_2024}. These models are supported by substantial neural evidence and underlying computational principles that describe the noisy accumulation of evidence leading to a choice \citep{shadlen_speed_2006, steinemann_direct_2024, bogacz_physics_2006, asadpour_rodents_2024, deakin_support_2024}.

\subsection{The Drift Diffusion Model (DDM)}
In this paper, we consider the DDM, a sequential sampling model which has a long history of describing decision making by a latent diffusion process that accumulates noisy sensory evidence over time until it reaches a decision bound, simultaneously explaining which choice is made and how long it takes to make it \citep{stone_models_1960, ratcliff_theory_1978, shadlen_neural_2001}. The model translates the latent diffusion process into observed behavior via parameters that together define the \emph{decision-making state}; elsewhere in the literature this is referred to as the decision-making strategy \citep{ashwood_mice_2022, roy_efficient_2018, roy_extracting_2021}. The DDM and its parameters are defined in Section~\ref{sec:ddm}.

\subsection{Decision Making Over Time}
The decision state is not constant; it changes over time with internal and external environments across timescales ranging from milliseconds to hours to days and beyond \citep{duffy_disentangling_2025, urai_choice_2019, cochrane_multiple_2023, shevinsky_interaction_2019, schumacher_neural_2023, schumacher_validation_2025}. Subjects adapt to these changing environments by adjusting their internal decision states, and these latent states in turn shape observed decision outcomes. Although standard measures of decision making often mask the temporal structure of these changes, recent approaches have begun to recover these time-varying dynamics by inferring parameter trajectories directly from behavioral data \citep{ashwood_mice_2022, gunawan_time-evolving_2022, kucharsky_hidden_2021, maggi_tracking_2024}. For example, subjects may alternate between discrete engaged stimulus-dependent states and disengaged stimulus-independent states \citep{ashwood_mice_2022, mohammadi_identifying_2025, duffy_disentangling_2025} or drift slowly between continuous states \citep{vloeberghs_bayesian_2025, roy_extracting_2021, roy_efficient_2018}, or a combination of both \citep{schumacher_neural_2023, schumacher_validation_2025}. As noted by \citet{urai_structure_2025} in the context of decision making, ``there is probably no such thing as stable behavior.''

Moreover, decisions are history dependent, influenced by sequential effects including the outcomes of previous choices \citep{frund_quantifying_2014, gao_sequential_2009}. Subjects use past experience to guide future choices, often integrating history information across multiple timescales \citep{danskin_exponential_2023, nguyen_optimizing_2019, busse_detection_2011}. This kind of history dependence can include effects of reinforcement learning, in which rewarded actions shape subsequent choices \citep{pedersen_drift_2017, fengler_beyond_2022}, as well as reward-independent sequential biases \citep{frund_quantifying_2014, gao_sequential_2009}. These forms of temporal structure are prevalent across perceptual decision making in humans, monkeys, and rodents \citep{urai_choice_2019, busse_detection_2011, trevino_response_2022, braun_adaptive_2018, gupta_trial-history_2024}. Humans may also exhibit forms of variability similar to those observed in rodents, differing only in the relative variability of different decision-making parameters \citep{nguyen_different_2022, shevinsky_interaction_2019}. Together, these demonstrate dynamic decision making across species, in which the past influences the present, producing trial-to-trial variability that is richly structured in time.

\subsection{Current Estimation Methods and Their Limitations}
Standard DDM fitting methods often implicitly assume that, in the absence of experimental manipulations, the decision-making state remains constant over time with fixed parameters, treating decision outcomes as independent and identically distributed ($iid$), a condition often violated in real data \citep{moens_recurrent_2018, urai_structure_2025, perquin_variance_2024}. When parameters are allowed to vary --- e.g., the seven-parameter DDM of \citet{ratcliff_modeling_1998} --- they do so as unstructured random variability around a stable mean, with trial order still ignored, and the possibility that systematic trends such as learning curves or circadian rhythms are linked to observable covariates is not considered \citep{ratcliff_parameter_2013}. Moreover, across-trial variability parameters are inherently difficult to estimate, showing poor recovery and weak test-retest reliability due to unidentifiability \citep{boehm_estimating_2018, lerche_model_2016, ratcliff_individual_2015, ratcliff_internal_2018, duffy_disentangling_2025, lerche_how_2017, van_ravenzwaaij_how_2009}.

Current inference methods compound these issues in four ways.
\begin{itemize}
    \item Point estimates do not always include corresponding uncertainties \citep{shinn_flexible_2020, boehm_estimating_2018}.
    \item When uncertainties are reported, they rely on asymptotic approximations that assume $iid$ conditions often not met in practice \citep{moens_recurrent_2018, wiecki_hddm_2013}.
    \item Quantifying uncertainty typically requires computationally expensive methods such as Bayesian MCMC or resampling \citep{boehm_estimating_2018, moens_recurrent_2018, wiecki_hddm_2013, fengler_hssm_2026}.
    \item The robustness of inference to unmodeled parameter variability and temporal correlations has not been addressed. Under misspecification, standard errors underestimate true uncertainty, resulting in overconfident and invalid inferences \citep{boehm_estimating_2018, urai_structure_2025, hato_levy_2025, schumacher_validation_2025}.
\end{itemize}

Together, these limitations constrain hypothesis testing and the detection of meaningful temporal changes in decision states, particularly in naturalistic settings where brain and behavior are in constant flux \citep{urai_structure_2025, ratcliff_parameter_2013}. As a result, even when rich temporal structure is present in long behavioral time series, experimenters lack a computationally efficient way to obtain valid uncertainties and hypothesis tests for time-varying DDM parameters.

\subsection{Proposed Solutions}
We consider the case of decision making over time, which introduces dependence and complicates inference. For example, parameters may trend within a session as satiety or fatigue accumulates, or drift between sessions as subjects learn. To address this, we explicitly allow for time dependence and model misspecification in order to conduct valid inference. We additionally propose modeling DDM parameters as functions of observed covariates (such as stimulus strength, trial number, or time of day) via link functions. This captures how external factors influence each decision and provides a principled mechanism for nonstationary behavior. Systematic parameter trends that would otherwise violate stationarity are explained by the covariate model, leaving residual variability that is plausibly stable over time. Stationarity is therefore a requirement on the joint process of covariates and outcomes, not on the observed decision outcomes alone. This distinction substantially broadens the scope of experiments to which the method applies.

We introduce a computationally efficient pseudo-maximum likelihood estimation (pseudo-MLE) method for estimating DDM parameters and covariate effects, together with analytical uncertainty estimates that remain valid under model misspecification and temporal dependence. The paper is organized as follows. Section~\ref{sec:methods} formalizes the model, estimator, and large-sample theory. Numbered equations indicate key expressions referenced throughout, while the others provide intermediate steps that can be skipped without loss of continuity. We implemented all described methods in Python except where other packages are explicitly cited. Section~\ref{sec:simulations} validates the method through simulations, and Section~\ref{sec:applications} demonstrates applications to rat decision making. Section~\ref{sec:practical-consideration} compares the proposed estimator with Bayesian MCMC, and Section~\ref{sec:conclusions} discusses implications, limitations, and future directions.
\end{document}
